\begin{textsample}{NEG Dim 2 – 2012\_c6 – Score -1.00 – 2012\_c6/enade\_2012\_tec\_gest\_rec\_hum\_c62ecb317d.txt}  \label{ex:f2_neg_185}
CAPÍTULO 3 ANÁLISE TÉCNICA DA PROVA Este capítulo tem por objetivo apresentar o desempenho dos estudantes concluintes de Tecnologia em Gestão de Recursos Humanos no ENADE/2012.
Para isso, foram calculadas as estatísticas básicas da prova em seu todo, bem como as estatísticas dos componentes relacionadas à Formação Geral, ao de Conhecimento Específico da Área e das questões discursivas isoladamente.
Nas tabelas, são apresentados o tamanho da população inscrita e de presentes, e as seguintes estatísticas das notas⁸ : média do desempenho na prova, erro padrão da média, desvio padrão, nota mínima, mediana e nota máxima.
Tais estatísticas contemplam o total de estudantes concluintes da Área de Tecnologia em Gestão de Recursos Humanos presentes à prova do ENADE 2012, tendo -se em vista as seguintes agregações : ( a ) as Grandes Regiões e o país como um todo ; ( b ) a Categoria Administrativa ; e ( c ) a Organização Acadêmica.
Em relação aos gráficos de distribuição de notas, o intervalo considerado foi de 10 unidades, aberto à esquerda e fechado à direita, com exceção do primeiro intervalo, [ 0 ; 10 ], fechado em ambos os extremos.
Para os gráficos de distribuição das notas das questões discursivas, foram consideradas mais duas categorias : questão em branco e nota zero. 3.1 ESTATÍSTICAS BÁSICAS DA PROVA 3.1.1 Estatísticas Básicas Gerais A Tabela 3.1 apresenta as estatísticas básicas da prova por grande Região.
A população total de inscritos foi de 48.586.
Destes, 37.611 estiveram presentes, sendo 22,6 \% o índice de não comparecimento.
A Região de maior abstenção foi a Centro-Oeste ( 31,7 \% ) e a de menor abstenção foi a Sul ( 14,9 \% ).
A média das notas da prova como um todo ( nas seções seguintes serão analisados os componentes de Formação Geral e de Conhecimento Específico ) foi 38,3, sendo que os alunos da região Centro-Oeste obtiveram a média mais baixa ( 35,7 ) e os da região Sul obtiveram a média mais alta ( 39,4 ).
As demais médias foram : 38,8 na região Sudeste, 37,9 na região Nordeste e 36,8 na região Norte.
O desvio padrão para o Brasil como um todo foi ⁸ Essas estatísticas e outras estão definidas no Capítulo 1. 11,9, sendo o maior desvio padrão encontrado na região Nordeste ( 12,1 ) e o menor na região Norte ( 11,0 ), indicando uma menor dispersão das notas desta última região.
A região que obteve a maior nota máxima foi a Sudeste ( 84,5 ), ao passo que a região que atingiu a menor nota máxima foi a Norte ( 64,2 ).
A mediana do Brasil como um todo foi 38,2, sendo a maior mediana obtida na região Sul ( 39,5 ) e a menor obtida na Centro-Oeste ( 35,5 ).
A nota mínima foi zero em todas as regiões sem exceção.
O comportamento das notas dos estudantes de todo o Brasil pode ser observado no Gráfico 3.1 que apresenta um histograma com a distribuição das mesmas.
Essa é uma distribuição unimodal com moda no intervalo ( 30 ; 40 ].
Apesar do coeficiente de assimetria da distribuição das notas ser positivo ( 0,05 ) este é pequeno e podemos considerar que a distribuição é aproximadamente simétrica.
As distribuições por Grande Região também apresentam assimetria positiva, concentração pouco maior do lado esquerdo do histograma e mais espalhada do lado direito, com exceção da região Sul onde a distribuição é simétrica ( 0,00 ).
Os Gráficos 3.2, 3.3 e 3.4 apresentam informações referentes à média da nota final dos Participantes, desagregando os resultados de acordo com, respectivamente, as Grandes Regiões do país, a Categoria Administrativa e a Organização Acadêmica.
Os gráficos apresentam o valor da média das notas como uma barra e os extremos do intervalo de confiança de 95 \% como linhas verticais unidas por uma linha horizontal na forma da letra H maiúscula.
Considerando -se o gráfico de notas segundo Grande Região ( Gráfico 3.2 ), observa -se que existe diferença estatisticamente significativa ao nível de 95 \% entre a maior média, obtida na região Sul ( 39,4 ) e a menor, obtida na região Centro-Oeste ( 35,7 ).
Levando -se em conta os estudantes por Categorias Administrativas ( Gráfico 3.3 ), observa -se que existe diferença estatisticamente significante entre as médias das notas das IES Públicas e Privadas.
Pode ser observado ainda que a média dos alunos de IES Públicas ( 41,7 ) é maior que a dos alunos de IES Privadas ( 38,3 ).
A diferença entre IES Públicas e Privadas ( 3,4 ) é menor do que a maior diferença entre regiões ( 3,7 ).
Tendo como foco o Gráfico 3.4, que apresenta as notas médias das provas segundo Organização Acadêmica, constata -se que existe diferença estatisticamente significativa ao nível de 95 \% entre a média das notas dos estudantes provenientes de Universidades e as médias dos alunos de Centros Universitários e Faculdades.
Entretanto, Centros Universitários e Faculdades não apresentam diferença estatisticamente significativa ao nível de 95 \% entre a média das notas.
A maior média foi obtida pelos estudantes de Faculdades ( 38,7 ), e a menor, pelos de Universidades ( 38,0 ). 3.1.2 Estatísticas Básicas no Componente de Formação Geral A Tabela 3.2 apresenta as estatísticas básicas em relação ao componente da prova que avalia a Formação Geral dos estudantes concluintes.
Os alunos de todo Brasil obtiveram desempenho médio de 38,6.
Quanto à variabilidade, o desvio padrão das notas dos estudantes do Brasil como um todo foi 15,2.
A maior média foi obtida na região Sul ( 39,8 ), e a menor, na região Norte ( 36,8 ).
As demais médias foram : 38,7 na região Sudeste, 37,7 na região Nordeste e 37,4 na região Centro-Oeste.
Já o maior desvio padrão foi obtido na região Centro-Oeste ( 15,5 ) e o menor na região Sudeste ( 15,1 ).
Os demais desvios padrões foram : 15,4 na região Nordeste, 15,3 na região Norte e 15,2 na região Sul.
A maior nota no componente de Formação Geral da prova do ENADE foi obtida por pelo menos um aluno nas regiões Sudeste e Centro-Oeste ( 92,0 ) enquanto que a menor nota máxima foi obtida na região Norte ( 76,5 ).
Nas outras regiões as notas máximas foram : 89,0 na região Sul e 87,5 na região Nordeste.
A mediana do Brasil como um todo foi 38,5, sendo a menor mediana encontrada na região Norte ( 36,5 ) e a maior encontrada na região Sul ( 40,0 ).
A nota mínima nesta parte foi zero em todas as regiões, sem exceção.
O Gráfico 3.5 propicia a avaliação do desempenho dos estudantes no componente de Formação Geral a partir do histograma da distribuição das notas correspondentes.
A distribuição é unimodal, com moda em ( 30 ; 40 ], mesmo intervalo modal da prova como um todo.
Nota -se, ainda, que no gráfico 3.5 as notas apresentam uma maior dispersão do que no Gráfico 3.1 ( distribuição das notas da prova ), confirmado pela comparação dos desvios padrões : 11,9 para a nota da prova como um todo e 15,2 para o componente de Formação Geral.
Para o componente de Formação Geral, o coeficiente de assimetria da distribuição das notas dos estudantes, como na prova como um todo, também é positivo e próximo de zero ( 0,02 ).
Por ser um coeficiente pequeno, a distribuição é praticamente simétrica.
Ao contrário das demais Grandes Regiões, os histogramas das regiões Norte e Sul possuem assimetria levemente negativa ( respectivamente -0,01 e -0,03 ).
Nos Gráficos 3.6, 3.7 e 3.8 são apresentadas as informações referentes ao desempenho dos Concluintes no componente de Formação Geral, em diferentes agregações : Grande Região do país, Categoria Administrativa e Organização Acadêmica.
Observa -se pelo Gráfico 3.6 que existe diferença estatisticamente significativa entre a maior ( 39,8 ) e a menor ( 36,8 ) média das notas no Componente de Formação Geral, respectivamente das regiões Sul e Norte.
Vemos que o intervalo de confiança mais largo é o da região Norte e o mais estreito é observado na região Sudeste.
Este fato está relacionado com o tamanho da população envolvida, menor na região Norte do que na região Sudeste.
No Gráfico 3.7, que representa as notas médias no Componente de Formação Geral segundo Categoria Administrativa do país, observa -se que existe diferença estatisticamente significativa entre as médias.
Assim como ocorreu para as notas da prova como um todo, para o componente de Formação Geral os concluintes de Tecnologia em Gestão de Recursos Humanos das IES Públicas ( 42,9 ) obtiveram uma média maior do que os das IES Privadas ( 38,6 ).
Considerando -se o tipo de Organização Acadêmica, nota -se, no Gráfico 3.8, uma diferença estatisticamente significativa entre a maior e a menor média.
Nas Universidades ( 38,8 ), a média é mais elevada do que em Centros Universitários ( 38,6 ) e Faculdades ( 38,2 ). 3.1.3 Estatísticas Básicas do Componente de Conhecimento Específico A Tabela 3.3 apresenta as estatísticas básicas referentes ao componente de Conhecimento Específico da área de Tecnologia em Gestão de Recursos Humanos.
A média do desempenho dos alunos do Brasil como um todo foi 38,2.
A maior média foi obtida na região Sul ( 39,2 ), e a menor, na região Centro-Oeste ( 35,1 ).
As demais médias foram : 38,8 na região Sudeste, 37,9 na região Nordeste e 36,8 na região Norte.
Quanto à variabilidade das notas, o desvio padrão do Brasil como um todo foi 13,1, sendo o maior desvio padrão observado na região Nordeste ( 13,3 ) e o menor na região Norte ( 11,9 ).
Os demais desvios foram : 13,1 da região Sul, 13,0 da região Sudeste e 12,9 da região Centro-Oeste.
A mediana das notas dos estudantes de todo o Brasil foi 38,0.
A maior mediana ocorreu na região Sul ( 39,5 ) e a menor na região Centro-Oeste ( 34,7 ).
As demais medianas foram 38,5 na região Sudeste, 37,4 na região Nordeste e 36,7 na região Norte.
A nota máxima do Brasil como um todo foi 86,6, sendo obtida por pelo menos um aluno da região Sudeste.
As demais notas máximas foram : 81,9 na região Sul, 81,1 na região Centro-Oeste, 78,0 na região Nordeste e 68,3 na região Norte.
A nota mínima foi zero em todas as regiões, sem exceção.
Assim como os Gráficos 3.1 e 3.5, o Gráfico 3.9, apresentado a seguir, proporciona uma avaliação do desempenho de concluintes em relação ao componente de Conhecimento Específico com um histograma da distribuição das notas correspondentes.
Esta também é uma distribuição unimodal, e novamente o grupo modal é o ( 30 ; 40 ].
O coeficiente de assimetria da distribuição das notas do componente de Conhecimento Específico é positivo e bem próximo de zero ( 0,07 ).
Nota -se pelo histograma ( Gráfico 3.9 ) que esta é uma distribuição aproximadamente simétrica.
Já na região Norte, o coeficiente de assimetria é negativo, igual a -0,06, evidenciando que a distribuição das notas dos estudantes dessa região, no componente de Conhecimento Específico, tem cauda mais pesada à esquerda.
Os Gráficos 3.10, 3.11 e 3.12 apresentam, respectivamente, uma comparação dos resultados em relação à Grande Região do país, à Categoria Administrativa e à Organização Acadêmica, agora levando em conta o desempenho de alunos no componente de Conhecimento Específico da prova.
Pelo Gráfico 3.10, observa -se que há diferença estatisticamente significativa entre a média das notas no Componente de Formação Geral, segundo Grande Região do país, se comparada à região Centro-Oeste com as demais regiões.
Sendo que a região Sul foi a que obteve média mais elevada ( 39,2 ), enquanto a região Centro-Oeste foi a região que obteve média mais baixa ( 35,1 ).
Quanto à Categoria Administrativa ( Gráfico 3.11 ), observa -se um comportamento semelhante àquele da parte de Formação Geral e à prova como um todo, ou seja, existe diferença estatisticamente significativa entre as médias das IES Públicas e Privadas.
As IES públicas de ensino também obtiveram a maior média ( 41,2 ).
Quanto ao Gráfico 3.12, observa -se, mais uma vez, que existe diferença estatisticamente significativa ao nível de 95 \% entre as notas, no componente de Conhecimento Específico, dos estudantes provenientes de Universidades e dos provenientes de Centros Universitários e Faculdades.
Sendo que a média dos Concluintes das Faculdades ( 38,8 ) foi maior do que de Centros Universitários ( 38,4 ) e de Universidades ( 37,6 ). 3.2 ANÁLISE DAS QUESTÕES OBJETIVAS 3.2.1 Componente de Formação Geral A Tabela 3.4 apresenta as estatísticas básicas relativas às oito questões objetivas do componente da prova que abrange a Formação Geral dos estudantes.
A média do Brasil foi 40,5.
A menor média foi encontrada na região Norte ( 38,2 ) e a maior na região Sul ( 42,2 ).
As demais médias foram 40,3 na região Sudeste, 39,9 na região Centro-Oeste e 39,8 na região Nordeste.
O desvio padrão do Brasil foi 19,4, sendo o maior desvio padrão encontrado na região Centro-Oeste ( 19,5 ) e o menor na região Norte ( 18,7 ).
Os demais desvios foram : 19,4 na região Nordeste, 19,3 na região Sudeste, 19,2 na região Sul.
As medianas ( 37,5 ) e as notas mínimas ( 0,0 ) foram iguais, nas questões objetivas do componente de Formação Geral, para todas as regiões.
As notas máximas ( 100,0 ) foram iguais para todas as regiões, com exceção da região Norte ( 87,5 ).
A Tabela 3.5 apresenta os índices de facilidade e discriminação ( ponto bisserial ) das questões objetivas do componente de Formação Geral.
Quanto ao índice de facilidade, foram usadas as seguintes cores para diferenciar o nível de dificuldade da questão : • Azul para as questões classificadas com índice muito fácil ( índice \&gt ; =0,86 ), verde para as questões classificadas com índice fácil ( 0,61 a 0,85 ), amarelo para as questões classificadas com médio ( 0,41 a 0,60 ), vermelho para as questões classificadas com difícil ( 0,16 a 0,40 ) e roxo para as questões classificadas com muito difícil ( \&lt ; =0,15 ).
Já quanto ao índice de discriminação, foram usadas as seguintes cores para qualificar a questão : • As questões classificadas com índice fraco receberam a cor roxa ( índice \&lt ; =0,19 ), as classificadas com médio receberam a cor vermelho ( 0,20 a 0,29 ), as classificadas com bom receberam a cor verde ( 0,30 a 0,39 ) e as classificadas com muito bom ( \&gt ; =0,40 ) receberam a cor azul.
As questões objetivas do componente de Formação Geral, segundo o índice de facilidade, foram assim avaliadas : das oito questões, três foram consideradas de dificuldade média, situando -se no intervalo entre 0,41 e 0,60 do índice de facilidade, ou seja, houve entre 41,0 \% e 60,0 \% de acertos ; as demais questões foram classificadas como difíceis, situando -se no intervalo entre 0,16 e 0,40.
Como já comentado, para análise das questões objetivas relativas à Formação Geral segundo o poder de discriminação, utilizou -se, o índice de discriminação ponto bisserial.
Nesta análise as questões foram assim avaliadas : cinco das oito questões apresentaram A Tabela 3.5 apresenta os índices de facilidade e discriminação ( ponto bisserial ) das questões objetivas do componente de Formação Geral.
Quanto ao índice de facilidade, foram usadas as seguintes cores para diferenciar o nível de dificuldade da questão :  Azul para as questões classificadas com índice muito fácil ( índice \&gt ; =0,86 ), verde para as questões classificadas com índice fácil ( 0,61 a 0,85 ), amarelo para as questões classificadas com médio ( 0,41 a 0,60 ), vermelho para as questões classificadas com difícil ( 0,16 a 0,40 ) e roxo para as questões classificadas com muito difícil ( \&lt ; =0,15 ).
Já quanto ao índice de discriminação, foram usadas as seguintes cores para qualificar a questão :  As questões classificadas com índice fraco receberam a cor roxa ( índice \&lt ; =0,19 ), as classificadas com médio receberam a cor vermelho ( 0,20 a 0,29 ), as classificadas com bom receberam a cor verde ( 0,30 a 0,39 ) e as classificadas com muito bom ( \&gt ; =0,40 ) receberam a cor azul.
As questões objetivas do componente de Formação Geral, segundo o índice de facilidade, foram assim avaliadas : das oito questões, três foram consideradas de dificuldade média, situando -se no intervalo entre 0,41 e 0,60 do índice de facilidade, ou seja, houve entre 41,0 \% e 60,0 \% de acertos ; as demais questões foram classificadas como difíceis, situando -se no intervalo entre 0,16 e 0,40.
Como já comentado, para análise das questões objetivas relativas à Formação Geral segundo o poder de discriminação, utilizou -se, o índice de discriminação ponto bisserial.
Nesta análise as questões foram assim avaliadas : cinco das oito questões apresentaram índices acima de 0,40 e, assim, foram classificadas com índice muito bom para esse grupo de alunos ; as outras três questões tiveram bom índice de discriminação, entre 0,30 e 0,39.
O índice de facilidade variou de 0,29 a 0,48, e o de discriminação, de 0,33 a 0,44.
Nem todas as questões com índices de discriminação muito bom, de números 1, 2, 3, 6 e 7, figuraram entre as mais fáceis desse conjunto : duas foram classificadas na categoria média ( questões 1 e 6 ) do índice de facilidade e três na categoria difícil ( questões 2, 3 e 7 ).
Em particular, a questão 3 foi a segunda mais difícil, com uma proporção de 0,37 de acertos.
A questão de número 4 apresentou índice de facilidade 0,46, ou seja, um quantitativo de 46,0 \% dos estudantes conseguiu resolvê -la, dentro do universo de participantes, e teve bom índice de discriminação ( 0,39 ).
A questão 8 obteve índice de discriminação médio, 0,37, e seu índice de facilidade foi difícil.
Já a questão de número 5, considerada difícil, obteve índice de discriminação 0,33, considerado bom.
O Gráfico 3.13, para exemplificar, analisa o comportamento da questão de número 7 de Formação Geral.
Trata -se da questão mais fácil e uma das que obteve o maior índice de discriminação dessa parte da prova.
Neste gráfico, cada uma das cinco curvas representa o percentual de respostas em determinada alternativa da questão, em função da nota dos estudantes nesta parte da prova ( Formação Geral/Múltipla Escolha ), antes de possíveis eliminações pelo critério do ponto bisserial.
A curva em verde corresponde à alternativa B, a correta para esta questão.
Assim, observa -se que entre os estudantes com notas mais baixas, nessa parte do exame, a situação mais frequente foi a escolha de uma das alternativas incorretas : a alternativa C ( em preto ) ou D ( em roxo ).
Na medida em que a nota aumenta, indicando desempenho melhor nesta parte da prova, aumenta concomitantemente a proporção de estudantes que selecionaram a alternativa correta B, atingindo 100 \% para os alunos com 8 acertos.
Essa análise permite verificar como a questão discriminou os grupos de desempenho, justificando o alto índice obtido na questão.
Os gráficos relativos às demais questões de Formação Geral constam do Anexo I. 3.2.2 Componente de Conhecimento Específico A Tabela 3.6 apresenta as estatísticas básicas em relação às questões objetivas do componente de Conhecimento Específico da prova, por Grande Região.
A média do Brasil deste componente foi de 42,5.
A menor média foi observada na região Centro-Oeste ( 39,4 ) e a maior na região Sul ( 43,7 ).
O desvio padrão de todo o Brasil foi 14,5, sendo o menor desvio padrão encontrado na região Norte ( 13,3 ) e o maior na região Nordeste ( 14,8 ).
A mediana de todo o Brasil foi 43,5, ocorrendo também nas regiões Nordeste, Sudeste e Sul do país.
Nas demais regiões a mediana foi 39,1.
A nota máxima da prova ( 100,0 ) foi obtida, nas questões objetivas do componente de Conhecimento Específico, por pelo menos um aluno da região Sudeste.
As notas máximas nas demais regiões foram : 91,3 nas regiões Sul e Centro-Oeste, 87,0 na região Nordeste e 73,9 na região Norte.
A nota mínima foi zero em todas as regiões do Brasil, sem exceção.
A Tabela 3.7 apresenta os índices de facilidade e discriminação ( ponto bisserial ) das questões objetivas do componente de Conhecimento Específico.
Para facilitar a diferenciação das questões usaram -se as mesmas cores da Tabela 3.5 para as diferentes classificações dos índices de facilidade e de discriminação.
A partir dos índices obtidos, pode -se concluir que a maioria das questões objetivas da prova foi considerada pelo menos difícil : das 27 questões, doze foram classificadas como difíceis e duas como muito difíceis.
Não houve questão classificada como muito fácil, ao passo que três foram tidas como fáceis, na faixa de 0,61 a 0,85 do índice de facilidade, e outras dez consideradas médias, entre 0,41 e 0,60.
Já quanto aos índices de discriminação das questões objetivas do componente de Conhecimento Específico da prova, tem -se como resultado a seguinte classificação : sete das 27 questões válidas foram consideradas como boas, enquanto três delas tiveram índice de discriminação muito bom.
Assim, para dez das 27 questões, os índices de discriminação foram bons ou muito bons.
Dentre as demais, treze delas foram classificadas como médias e outras quatro como fracas, sendo dezessete, por conseguinte, a quantidade de questões nos dois patamares mais baixos de discriminação.
Constata -se, assim, que a prova – no que se refere ao componente de Conhecimento Específico – possuía capacidade fraca de discriminar entre aqueles que dominam ou não o conteúdo.
Dentre as questões que alcançaram os maiores índices de discriminação, as de números 14, 21 e 35 classificadas com índice muito bom, situando -se no intervalo de 0,40 a 0,45 do índice, duas delas ( questões 14 e 35 ) foram classificadas na categoria fácil e uma ( 21 ) como média, quanto ao índice de facilidade.
As questões de número 23 e 24 foram as mais difíceis dentre as 27 questões específicas, com baixo índice de facilidade, apenas 12,0 \% de acertos.
Essas questões apresentaram poderes discriminatórios igualmente baixos, respectivamente 0,10 e 0,02, o que comprova terem sido estas as mais difíceis para os estudantes.
Destaca -se, também, a questão 20, com índice de facilidade 0,24, o que, em termos percentuais, corresponde a 24,0 \% de estudantes que responderam acertadamente, obtendo, ainda, 0,14 de índice de discriminação.
Tais questões foram, portanto, três das mais difíceis da prova.
Pelo critério do ponto bisserial, elas foram eliminadas do cômputo da nota final.
Além dessas três, a outra questão com índice fraco de discriminação, questão 13 também não foi computada.
A título de exemplo das análises do comportamento das questões objetivas, o Gráfico 3.14 analisa a questão 21 do componente de Conhecimento Específico.
Esta foi apenas a sétima questão mais fácil da prova, apresentando índice de facilidade 0,52, ou seja, 52,0 \% dos estudantes assinalaram acertadamente a opção A, correspondente ao gabarito.
Já seu índice de discriminação foi o maior ( 0,45 ), classificado como muito bom.
Neste gráfico, cada uma das cinco curvas representa o percentual de respostas em determinada alternativa da questão 21, em função da nota dos estudantes nesta parte da prova, antes de possíveis eliminações de questões pelo critério do ponto bisserial.
A alternativa correta A, representada no gráfico pela curva em azul, foi escolhida em maiores proporções pelos alunos com desempenho melhor nesta parte da prova.
Já as alternativas incorretas, também denominadas distratores, foram selecionadas principalmente por Aqueles com nota zero, na sua quase totalidade deixaram esta questão em branco ou marcaram mais de uma alternativa, comportamento considerado inválido.
A proporção de alunos que selecionou a resposta correta A aumenta gradativamente, chegando a atingir 100 \% para os alunos com 21 acertos ou mais, enquanto a proporção dos que escolheram alternativas incorretas decai, a partir da primeira nota não nula, como função da nota nesta parte da prova.
Os gráficos relativos às demais questões do componente de Conhecimento Específico constam do Anexo I. 3.3 ANÁLISE DAS QUESTÕES DISCURSIVAS 3.3.1 Componente de Formação Geral As análises dos resultados de desempenho dos estudantes de Tecnologia em Gestão de Recursos Humanos nas duas questões discursivas relativas à Formação Geral encontram -se na Tabela 3.8 e no Gráfico 3.15.
Na tabela 3.8 observa -se que as notas médias foram mais baixas nesse conjunto de questões do que no das objetivas.
Os estudantes de todo o Brasil obtiveram, em Formação Geral, média 40,5 nas questões objetivas e 35,7 nas questões discursivas.
A mediana dos alunos de todo o Brasil nas questões discursivas do componente de Formação Geral foi a mesma ( 37,5 ) observada para as questões objetivas.
Pode -se, também, notar um leve aumento do desvio padrão de 19,4, nas questões objetivas do componente de Formação Geral dos alunos de todo o Brasil, para 20,0 nas questões discursivas do mesmo componente.
As regiões Norte, Sudeste e Sul obtiveram a mesma mediana do país como um todo ( 37,5 ), enquanto as demais regiões tiveram mediana igual a 35,0.
A nota máxima ( 90,0 ) foi obtida, nas questões discursivas do componente de Formação Geral, por pelo menos um aluno da região Sudeste, enquanto o menor valor ( 80,0 ) foi obtido na região Norte.
A nota mínima foi zero em todas as regiões do Brasil.
O Gráfico 3.15 representa a distribuição das notas nas questões discursivas no componente de Formação Geral.
A moda desta distribuição ocorre no intervalo, ( 40 ; 50 ].
Pode -se também observar uma moda secundária no primeiro intervalo, [ 0 ; 10 ].
A distribuição possui assimetria à esquerda, coeficiente de assimetria -0,30.
Pelo gráfico pode -se observar que esse valor é influenciado pela concentração de notas no intervalo [ 0 ; 10 ].
Em todas as regiões o coeficiente de assimetria também é negativo, o maior em módulo, -0,33, nas regiões Sudeste e Sul, e o menor nas regiões Centro-Oeste e Nordeste ( -0,22 ).
Na sequência, os resultados verificados para cada uma das questões discursivas de Formação Geral serão apresentados, estabelecendo -se relações com os conteúdos abordados em cada uma delas.
Os comentários da Banca de docentes corretores a respeito do observado na correção das respostas dos estudantes, suas impressões e conclusões serão apresentados junto à análise de cada questão.
Cumpre esclarecer que, tendo em vista que as questões discursivas de Formação Geral são padronizadas, ou seja, constam de todas as provas, os comentários da Banca são os mesmos para todas as carreiras acadêmicas, sendo direcionados a todos os estudantes que participaram do ENADE/2012.
A seguir, serão analisados os desempenhos da Área de Tecnologia em Gestão de Recursos Humanos nas duas questões discursivas de Formação Geral do ENADE/2012, comparando os resultados obtidos com comentários para cada questão. 3.3.1.1 Análise da Questão Discursiva 1 do Componente de Formação Geral Os dados de Tecnologia em Gestão de Recursos Humanos, obtidos a partir das respostas à questão 1, encontram -se na Tabela 3.9 e no Gráfico 3.16.
Nessa questão – de pior desempenho dentre as duas de Formação Geral – os alunos de todo Brasil tiveram média, 28,2.
A maior média para a questão 1 foi obtida nas regiões Sudeste e Sul ( 29,0 ), e a menor, na região Centro-Oeste ( 25,9 ).
Quanto à variabilidade das notas, o desvio padrão de todo o Brasil foi 23,9.
O menor desvio padrão foi obtido na região Sul ( 23,6 ) e o maior desvio padrão foi obtido na região Norte ( 24,3 ).
A mediana do país como um todo ( 30,0 ) foi a mesma nas regiões Sudeste e Sul.
Nas demais regiões as medianas foram : 25,0 na região Nordeste e 20,0 nas regiões Norte e Centro-Oeste.
As notas máximas da questão discursiva 1 foram as mesmas para quase todas as regiões do Brasil ( 95,0 ), com exceção da região Norte, cuja máxima foi 90,0.
Além disso, a nota mínima ( 0,0 ) foi obtida em todas as regiões do país.
O Gráfico 3.16 mostra a distribuição das notas na questão discursiva 1 do componente de Formação Geral.
Observa -se que a maior frequência corresponde aos alunos que deixaram a questão em branco, a classe modal..
Considerando apenas os estudantes que tentaram resolver a questão, a moda ocorre no intervalo ( 20 ; 30 ].
A distribuição é assimétrica à direita, com coeficiente de assimetria positivo ( 0,34 ).
Em todas as regiões do Brasil este coeficiente se mantém positivo, variando entre 0,30 na região Sul e 0,48 na região Norte. 3.3.1.2 Comentários sobre a correção das respostas à Questão Discursiva 1 A questão aborda um tema bastante polêmico o qual vem sendo amplamente discutido, a conciliação entre o desenvolvimento econômico e a preservação dos recursos naturais.
Busca -se observar a relação entre o estímulo governamental à compra de veículos novos e os impactos ambientais decorrentes desta política pública.
Sua compreensão implica no conhecimento de uma temática atual e na reflexão critica a respeito das decisões governamentais em sua correlação com o cenário econômico internacional e os problemas e soluções de ordem ambiental.
Embora o comando da questão solicitasse explicitamente que a resposta fosse fornecida em um texto único ( “ redija um texto dissertativo sobre... ” ), foi significativamente alta a proporção de respondentes que atendeu a cada item ( a, b e c ) do enunciado de maneira isolada, como partes estanques da mesma resposta.
No que se segue encontra -se uma avaliação da correção de cada um dos três itens. a - “ conceito de desenvolvimento sustentável ” A variedade de respostas a esse item confirmou a importância do tema para avaliar a formação geral dos estudantes e cursos, já que ( i ) é atual ; ( ii ) foi abordado nos meios de comunicação ( especialmente no ano de 2012, em que o Brasil sediou a “ Rio+20 ” ) ; ( iii ) é relevante não apenas no Brasil como em todos os principais fóruns internacionais ; e ( iv ) é comum a diversas áreas de conhecimento.
Outros dois aspectos positivos merecem destaque : ( i ) a grande incidência de respostas que apontaram a relação entre gerações como aspecto essencial da ideia de sustentabilidade socioambiental – no mais, as respostas apresentaram pelo menos a tentativa de indicar uma definição de desenvolvimento sustentável ou a abordagem positiva dessa mesma noção, evidenciado compreensão do enunciado ; e ( ii ) o exercício crítico que a questão propunha, no sentido de exigir dos respondentes a compreensão das contradições de políticas públicas e a apresentação de propostas de soluções para conciliar o incentivo do crescimento econômico e a promoção da sustentabilidade socioambiental, a partir de uma perspectiva transformadora da realidade.
Apesar dos aspectos positivos apontados no parágrafo anterior, a maioria dos estudantes não alcançou o máximo da pontuação.
De forma geral, trataram o desenvolvimento sustentável como sinônimo de não agressão ambiental.
Alguns confundiram desenvolvimento sustentável com desenvolvimento econômico, e outros faziam referência apenas à preocupação com a conservação dos recursos mencionando a preocupação com gerações futuras.
Deve -se, ainda, destacar o significativo número de alunos que apresentou uma definição aparentemente copiada do enunciado de outra questão da prova, como por exemplo, a questão objetiva 31 da prova de Administração. b - “ conflito entre o estímulo à compra de veículos automotores e a promoção da sustentabilidade ” Esse foi o item identificado como aquele no qual as respostas dos estudantes avaliados mais se aproximaram da pontuação máxima prevista no padrão de resposta.
Ou seja, a ampla maioria dos respondentes foi bem-sucedida em apontar a relação entre o aumento da quantidade de veículos automotores em circulação e os danos socioambientais que isso acarreta no cenário urbano, devido ao aumento da emissão de gases poluentes.
No item b, as respostas corretas se concentraram em críticas ao incentivo oferecido pelo governo, referências ao problema dos congestionamentos e ao aumento da poluição.
Contudo, a maioria das respostas não explicitava a correlação entre o aumento da circulação de veículos automotores e o agravamento dos problemas ambientais.
Um número expressivo de alunos apresentou uma resposta em que era dado destaque à política de redução do IPI como um mecanismo para renovação da frota de automóveis ou para a aquisição do transporte próprio e da consequente redução da dependência ao transporte coletivo.
Também foram notados alguns tipos de padrões recorrentes nas tentativas de resposta a este item.
Dois desses padrões se destacaram : ( i ) as declarações genéricas de condenação do “ consumismo ” de uma “ nova classe média ” ou da busca pelo “ status ” que a posse de um automóvel particular proporcionaria, como se fosse possível atribuir a essa ordem de coisas a causa principal do aumento dos congestionamentos nas grandes cidades ; e ( ii ) a repetição da avaliação, presente no enunciado da questão, que aponta de forma superficial a contradição entre o estímulo à demanda por automóveis particulares e a promoção da sustentabilidade. c - “ ações de fomento ao transporte urbano sustentável no Brasil ” Este foi o item que ofereceu maior facilidade de aplicação do padrão de resposta, mas foram poucos os casos de pontuação integral.
Dentre as soluções mais destacadas pelos alunos apareceram : ( i ) o incentivo ao uso de energias alternativas ; e ( ii ) o estímulo ao transporte coletivo ( metrô, trem e ônibus ), além do uso de bicicletas.
Importante destacar o fato de que muitos apresentavam respostas com referências à necessidade de incentivo ao transporte urbano, nas quais este termo era utilizado como sinônimo de transporte coletivo.
A avaliação da questão 1 trouxe à tona problemas estruturais da educação básica, que se refletem no ensino superior.
Foi muito grande o número de provas em que o aluno demonstrou não compreender o enunciado da questão, apresentando respostas que fugiam ao tema proposto, mesmo quando bem redigidas.
Uma resposta recorrente fazia referência às consequências da ampla oferta de crédito à população e aos problemas relacionados à inadimplência, situação que, embora, verdadeira, não estava relacionada à questão proposta.
Houve, ainda, os casos em que as respostas apresentadas não só se afastaram do tema como não reuniam qualquer tipo de argumentação coerente.
E, por fim, destaca -se a existência de respostas incoerentes e desconexas.
Outro aspecto geral diz respeito ao grande número de provas em branco ou em que o aluno apresentou a resposta a outro item no espaço da questão 1.
Considerando -se o objetivo geral deste exame, avaliar a qualidade dos cursos de graduação oferecidos no país e contribuir para a melhoria da qualidade do ensino superior, destaca -se a imensa responsabilidade do processo que envolve a elaboração deste diagnóstico.
Nesse sentido, é urgente alertar para a ocorrência de um número significativo de respostas que se afastam de um padrão mediano de qualidade esperado de um aluno concluinte em um curso de graduação, especialmente em se considerando o fato de que a questão proposta faz referência a tema intensamente explorado pela mídia. 3.3.1.3 Análise da Questão Discursiva 2 do Componente de Formação Geral A Tabela 3.10 mostra que o desempenho dos estudantes na questão 2 ( média 43,1 ) foi tipicamente superior ao obtido na questão de número 1 ( média 28,2 ).
A região Sudeste foi aquela onde a média, nessa questão, foi maior ( 43,7 ), e a de menor média foi a região Centro-Oeste ( 41,0 ).
Quanto à variabilidade das notas, o desvio padrão de todo o Brasil foi 23,4, inferior ao obtido na questão de número 1 ( 23,9 ).
O maior desvio nessa questão foi obtido nas regiões Nordeste e Centro-Oeste ( 24,4 ), enquanto o menor foi obtido na região Sul ( 22,8 ).
A mediana do país como um todo ( 50,0 ) só não foi a mesma nas regiões Norte e Centro-Oeste, que tiveram mediana 45,0.
A maior nota máxima ( 95,0 ) foi obtida por alunos das regiões Nordeste e Sul, enquanto a menor nota máxima ( 85,0 ) foi obtida na região Norte.
A nota mínima foi zero em todas as regiões do Brasil, sem exceção.
O Gráfico 3.17 mostra a distribuição das notas na questão discursiva 2 do componente de Formação Geral.
Neste gráfico observa -se a razoável quantidade de alunos que deixaram a questão 2 em branco ( em torno de 15 \% ).
A moda da distribuição está no intervalo ( 40 ; 50 ].
O coeficiente de assimetria das notas nesta questão é negativo ( -0,61 ), indicando que a distribuição possui assimetria à esquerda.
O mesmo ocorre para todas as regiões do país, com coeficientes variando entre -0,64 na região Sudeste e -0,51 na região Centro-Oeste. 3.3.1.4 Comentários sobre a correção das respostas à Questão Discursiva 2 O conteúdo da questão era bastante adequado às Diretrizes do INEP para o ENADE, além de atual e de abrangente.
Porém, exigia -se a articulação de muitas variáveis.
Buscando aferir, a partir da apresentação de uma charge e de um texto introdutório, as competências e habilidades relacionadas à leitura, à interpretação, à análise crítica, à capacidade argumentativa, à elaboração de sínteses e de avaliação prospectiva sobre tema da atualidade, a questão tratava da violência e suas repercussões no contexto atual.
O comando da questão solicitava, a partir da definição de violência da OMS e de duas imagens, que o aluno elaborasse um texto, de caráter dissertativo, sobre a “ violência na atualidade ”.
Embora tal solicitação fosse explícita ( “ redija um texto dissertativo sobre... ” ), como na questão 1, foi alta a proporção de alunos que respondeu a cada item do enunciado isoladamente, como partes independentes.
Verificaram -se poucas respostas em branco, como também foram poucas as respostas contendo protestos contra o ENADE, muito frequentes em anos anteriores, o que significa um avanço para a afirmação do ENADE como um exame aceito socialmente. a -"tecnologia e violência"Tanto as charges quanto o texto apresentado no enunciado induziram os alunos a focarem suas respostas nos meios de comunicação.
Além disso, algumas respostas estabeleciam a tecnologia, em si, como causa da violência tanto nas escolas ( em particular ), quanto na sociedade ( no geral ).
Muitos respondentes tiveram dificuldade para fazer a relação entre tecnologia, violência e a escola, de modo adequado, centrando -se na violência nas escolas.
Além disso, a grande maioria das respostas girava em torno da exposição da violência através da TV, internet e jogos digitais, sendo que \textbf{predominaram} as abordagens com ênfase em programas e filmes veiculados pela televisão.
Poucos foram os estudantes que estabeleceram relação entre o contexto socioeconômico e a violência. b -"causas e consequências da violência na escola"Embora a solicitação desse item fosse objetiva e atendida por grande parte dos alunos, inúmeras vezes, observou -se confusão entre o significado de causa e de consequência.
Em outras inúmeras vezes, as respostas apontavam apenas para uma das classes, causas ou consequências.
Foram frequentes as respostas que atribuíam o aumento da violência à ineficiência do poder público, apontando a falta de investimentos, o baixo salário dos professores, a ausência de profissionais qualificados como causas.
Foram frequentes também os que apontavam a violência doméstica, o bullying, e a transferência do imaginário dos filmes e jogos para a realidade, como causas da violência. c -"proposta de solução para o problema da violência na escola"Sem dúvida, há uma grande diversidade de abordagens possíveis para apresentação de uma proposta de solução para o problema da violência, o que o padrão de resposta buscou contemplar.
No entanto, houve grande número de respostas baseadas em opiniões pessoais dos estudantes e vinculadas ao senso comum.
As respostas giraram em torno da necessidade de mais investimentos na educação por parte do poder público e seu rebatimento no maior aparelhamento das escolas, salário dos professores etc.
Além disso, foram muitas as respostas que trataram da necessidade de maior conscientização dos pais, alunos e professores com a realização de um trabalho conjunto que envolvesse todos os segmentos da sociedade.
De um modo geral, os estudantes demonstraram dificuldades em relação ao nível de abrangência exigido, falta de conteúdo, inabilidade na elaboração de texto escrito.
Da grande quantidade de provas corrigidas, o que ficou evidente, foi a dificuldade de estabelecer as conexões requeridas na questão 2.
A violência nas escolas foi bastante associada ao reflexo do meio onde a criança e o adolescente mantém os seus vínculos afetivos.
Surgiram muitas respostas revelando críticas à intolerância no núcleo familiar, à jornada de trabalho dos pais, ao distanciamento da mãe nas atividades domésticas etc.
Em boa parte das respostas, os estudantes lançaram a causa da violência escolar exclusivamente na ausência de participação familiar na vida dos filhos, às vezes criando uma cisão com a proposta da questão, já que abandonavam totalmente a temática da tecnologia em suas reflexões.
Além da ausência dos pais no cotidiano da escola, divórcios e atos de violência contra crianças foram considerados fatores que delegariam para a escola novas funções.
Também foram mencionadas a má formação de professores e a insatisfação com o ambiente escolar.
Em inúmeros casos, ficou evidente o anseio de se fazer da escola um lugar de resgate tanto da criança e do adolescente quanto da família, compreendendo -se a escola em sentido social e assistencial.
Segundo este amplo grupo de concluintes, a escola deveria oferecer -se como centro de cultura e lazer, como espaço para o acompanhamento psicológico e serviço social às comunidades diretamente envolvidas.
Ressalte -se o constante aparecimento das sugestões de horário integral, melhores salários para professores e iniciativas para melhoria na formação de professores.
As propostas de solução para acabar com a violência na escola se detiveram em aspectos muito gerais e poucos alunos detalharam uma proposta de intervenção, apresentando ações concretas.
Dentre as propostas de solução surgiram : educação familiar, ensino religioso, maior segurança nas escolas, com utilização de aparelhos para detecção de armas, e o retorno do papel do professor como um profissional respeitado pelos poderes municipais e estaduais.
Um percentual significativo de provas apresentou resposta simplista, com ênfase em determinados pontos como : religião ( sem demonstração do conceito de laicidade ), políticas sociais ( sem demonstração do domínio de conceitos que envolvem Estado, Governo e Politicas Públicas ) etc.
Foi grande a incidência de respostas que se referiam às leis, à vigilância e maior controle dos pais e professores ; no entanto, não demonstram entender a noção de legalidade para os diferentes segmentos sociais e como os mesmos se encontram submetidos a diferentes formas de poder.
Alguns alunos, ainda que raros, trataram a questão com a profundidade esperada de concluintes de cursos universitários.
Por outro lado, foram vários os textos que apresentaram uma argumentação simplória e desconexa, com forte visão paternalista sobre o papel do governo e do professor, como se tais instâncias pudessem intervir na educação familiar.
Enquanto alguns conseguiram desenvolver o tema proposto com elegância e desenvoltura, encadeando o tema de forma harmônica, com excelente nível de linguagem escrita e correção gramatical, outros sequer conseguiram conectar as ideias e construir um texto razoável, utilizando, inclusive, um linguajar mais propício à linguagem falada do que à escrita.
Cabe também ressaltar a quantidade de respostas carregadas de conteúdo com juízo de valor discriminatório ou extremamente conservador, do tipo : “ mulher não deve trabalhar fora ” ; “ amar e temer a Deus é a solução ” ; obrigatoriedade do ensino religioso e de uma doutrina religiosa ; imposição de uma educação do tipo militar nas escolas ; censura nos meios de comunicação etc.
Outras soluções apontadas para resolver o problema da violência nas escolas estavam associadas à, por exemplo : redução da maioridade penal, defesa da pena de morte, milícias armadas na porta das escolas, punições como o retorno da palmatória.
O Estatuto da Criança e do Adolescente foi uma peça legal criticada, e recebeu, em diversas respostas, o título de responsável pelo sentimento de impunidade de adolescentes infratores.
Contudo, algumas dessas proposições, desde que bem contextualizadas e formuladas com boa argumentação, foram pontuadas de forma a respeitar o direito democrático de pluralidade opinativa.
É imperioso ressaltar que a falta de inteligibilidade segue representando um aspecto merecedor de providências na formação universitária, uma vez que exprime uma habilidade importante para qualquer profissional de nível superior.
Além de terem sido poucas respostas boas e claras, foi possível identificar a deficiência de argumentação, além dos inúmeros problemas gramaticais ( ortografia, concordância, pontuação, regência, etc. ).
Salienta -se, nesse sentido, o uso indiscriminado de letras maiúsculas e minúsculas numa mesma frase, sem obedecer ao critério formal da língua portuguesa. 3.3.2 Componente de Conhecimento Específico Na parte da prova relativa às questões discursivas no componente de Conhecimento Específico ( Tabela 3.11 ), observa -se que a média foi mais baixa do que para as questões discursivas do componente de Formação Geral.
Enquanto no componente de Formação Geral a média para estudantes de Tecnologia em Gestão de Recursos Humanos de todo o Brasil foi 35,7, na parte de Conhecimento Específico a média foi 13,6.
A maior média deste componente foi obtida pelos estudantes da região Sudeste ( 14,5 ), e a menor, pelos da região Centro-Oeste ( 10,9 ).
Quanto à variabilidade das notas, o desvio padrão de todo o Brasil foi 12,8.
O maior desvio padrão foi encontrado nas regiões Nordeste e Sudeste ( 13,1 ), e o menor, na região Centro-Oeste ( 11,7 ).
A maior nota máxima, nas questões discursivas do componente de Conhecimento Específico, foi obtida na região Sudeste ( 73,3 ) enquanto a menor nota máxima foi encontrada na região Norte ( 53,3 ).
Além disso, a mesma nota mínima ( 0,0 ) foi obtida por alunos de todas as regiões do Brasil sem exceção.
A mediana do Brasil como um todo foi 10,0, também encontrada nas regiões Norte e Nordeste.
As demais medianas foram : 13,3 nas regiões Sudeste e Sul e 6,7 na região Centro-Oeste.
O Gráfico 3.18 representa a distribuição das notas nas questões discursivas no componente de Conhecimento Específico.
Essa distribuição tem moda no intervalo de [ 0 ; 10 ].
Nota -se que a partir desse intervalo há concentração decrescente de notas nos intervalos seguintes, também sendo relevante não terem sido registradas notas acima de 80,0.
O coeficiente de assimetria é positivo ( 0,79 ).
Este coeficiente também é positivo e próximo de 1,00 em todas as regiões do país ( entre 0,71 na região Sudeste e 1,05 na região Centro-Oeste ). 3.3.2.1 Análise da Questão Discursiva 3 do Componente de Conhecimento Específico Na questão 3, cujos resultados aferidos encontram -se descritos na Tabela 3.12, a média dos estudantes de todo o Brasil foi 15,1.
A menor média nessa questão foi obtida pelos alunos da região Centro-Oeste ( 12,8 ), enquanto a maior média foi obtida na região Sudeste ( 15,7 ).
Quanto à variabilidade das notas, o desvio padrão de todo o Brasil foi 15,6.
O maior desvio padrão foi obtido na região Nordeste ( 16,2 ), enquanto o menor foi obtido na região Centro-Oeste ( 14,9 ).
A nota máxima, 100,0 pontos, foi alcançada por pelo menos um aluno da região Sudeste.
A região onde ocorreu a menor nota máxima foi a Norte ( 60,0 ).
A mediana do Brasil como um todo foi 10,0, e foi a mesma em todas as regiões.
A mesma nota mínima ( 0,0 ) foi obtida em todas as regiões do Brasil.
O Gráfico 3.19 mostra a distribuição das notas na questão discursiva 3, do componente de Conhecimento Específico.
Essa distribuição é bimodal, com modas nas questões em branco e no intervalo ( 20 ; 30 ] correspondendo a um máximo local.
A distribuição das notas dos alunos que resolveram a questão 3 possui assimetria à direita, com coeficiente de assimetria 0,89.
Em todas as regiões do país o coeficiente de assimetria também é positivo ( entre 0,80 na região Sul e 1,07 na região Centro-Oeste ). 3.3.2.2 Comentários sobre as respostas à Questão Discursiva 3 O texto de contextualização da questão era bem elaborado e muito atual, alinhado com a problemática vivenciada nas organizações, permitindo avaliar o (des)conhecimento dos alunos sobre o assunto.
O comando da questão solicitava “ Considerando o contexto acima, faça o que se pede ” seguido de três subitens.
Entretanto, foi verificado que os alunos tiveram dificuldade em responder aos três subitens separadamente.
Além de responderem em um único texto, muitos se detiveram apenas ao solicitado no item a.
No item a, o desconhecimento da palavra gap dificultou a interpretação da questão.
Além disso, como a pergunta direcionou para o processo de treinamento e não para a avaliação do desempenho, a maioria respondeu que o processo de treinamento desenvolve o profissional melhorando os resultados da empresa, sem detalhes.
Para os itens b e c as respostas foram genéricas em relação ao modelo de gestão por competências, embora os comandos fossem claros e bem elaborados. 3.3.2.3 Análise da Questão Discursiva 4 do Componente de Conhecimento Específico A Tabela 3.13 contém as informações relativas à questão 4 do conjunto de questões do componente de Conhecimento Específico.
O desempenho dos estudantes de todo o Brasil nesta questão foi tipicamente inferior ao desempenho na questão de número 3.
A média geral do Brasil foi 9,7, sendo a menor média registrada na região Norte ( 7,3 ) e a maior na região Sudeste ( 10,6 ).
A maior nota máxima ( 100,0 ) foi atingida na região Sudeste, enquanto a menor nota máxima ( 50,0 ) foi atingida na região Norte.
A mediana e a nota mínima foram zero, e obtidas em todas as regiões do Brasil, sem exceção, ou seja, pelo menos metade dos alunos tirou nota zero nesta questão.
O Gráfico 3.20, representa a distribuição de notas na questão discursiva 4, no componente de Conhecimento Específico.
Como nas outras questões discursivas, o número de estudantes que deixou a questão em branco foi muito grande, correspondendo à moda da distribuição.
Dentre as notas dos que resolveram a questão 4, as maiores frequências são das menores notas.
A assimetria da distribuição confirma o valor baixo da média, que foi 9,7, já que apresenta concentração em notas pequenas e uma cauda longa do lado direito da distribuição.
Tanto na distribuição como um todo ( 1,59 ) quanto na análise dos resultados dos alunos por região ( entre 1,47 na região Sudeste e 1,90 na região Centro-Oeste ), os coeficientes de assimetria são sempre positivos e com valor alto, sempre acima de 1. 3.3.2.4 Comentários sobre as respostas à Questão Discursiva 4 Apesar de o enunciado ser bem contextualizado e próprio ao momento atual da gestão de recursos humanos, a questão exigiu a integração de conteúdos diversos ; o que, levando -se em consideração a extensão de toda a prova, talvez tenha dificultado a redação de respostas boas e completas.
Ao longo da correção, ficou claro que a quase totalidade dos estudantes respondeu aos aspectos mais evidentes da questão, não discorrendo sobre os que, a princípio, pareciam ser menos relevantes.
Apesar da propriedade, atualidade e pertinência do conteúdo avaliado pela questão, o desempenho dos estudantes evidenciou despreparo e falta de conhecimento.
Ficou clara a dificuldade de interpretação e de raciocínio lógico, além de uma frágil e imprecisa construção de conhecimento específico.
Tal resultado mostra um comprometimento significativo para a prática desses futuros profissionais.
Quase sempre, as respostas apresentavam conteúdos memorizados e não apreendidos e aprendidos e, portanto, sem argumentação lógica própria.
Além de o resultado espelhar uma precária formação profissional, também foi possível observar a crônica e significativa falta de preparo básico e elementar da maioria dos estudantes. 3.3.2.5 Análise da Questão Discursiva 5 do Componente de Conhecimento Específico A Tabela 3.14 contém as informações relativas à questão 5 do conjunto do componente de Conhecimento Específico.
O desempenho dos estudantes nessa questão foi superior ao das questões 3 e 4.
A nota média dos estudantes de todo o Brasil foi 16,1.
A maior média foi registrada na região Sudeste ( 17,1 ), enquanto a menor média foi registrada na região Centro-Oeste ( 12,5 ).
Quanto à variabilidade das notas, o desvio padrão dos alunos do Brasil, como um todo, foi 19,4.
Enquanto o maior desvio foi encontrado na região Norte ( 19,9 ), o menor foi encontrado na região Centro-Oeste ( 17,7 ).
A maior nota máxima ( 100,0 ) foi alcançada na região Sudeste, enquanto a menor nota máxima ( 80,0 ) foi registrada nas regiões Norte e Centro-Oeste.
A mediana foi 10,0 para a maioria das regiões, exceto para as regiões Nordeste e Centro-Oeste, que tiveram mediana igual a zero.
A nota mínima foi zero para todas as regiões do Brasil, ou seja, pelo menos metade dos estudantes das regiões Nordeste e Centro-Oeste tirou nota zero nesta questão.
O Gráfico 3.21 representa a distribuição das notas da questão discursiva 5 do componente de Conhecimento Específico.
Destaca -se, como nas demais questões discursivas, o grande número de estudantes que deixaram esta questão em branco, correspondendo à moda. correspondendo à moda.
Entre as três questões discursivas do Componente de Conhecimento Específico, esta foi a que apresentou um maior espalhamento das notas daqueles que resolveram a questão.
Uma moda secundária ocorre no intervalo ( 10 ; 20 ].
O coeficiente de assimetria para todos os alunos é 0,96 e para todas as regiões se mantém positivo, variando entre 0,88 nas regiões Sudeste e Sul, e 1,29 na região Centro-Oeste.
Essa assimetria positiva indica que o gráfico possui maior concentração das notas do lado esquerdo do gráfico e uma cauda maior do lado direito. 3.3.2.6 Comentários sobre as respostas à Questão Discursiva 5 Tanto quanto a questão anterior — discursiva 4 —, essa questão apresenta um contexto bastante próprio ao momento atual da gestão de recursos humanos e exigiu a integração de assuntos diversos e interdependentes.
Consistente e precisa, a questão foi preparada de modo a possibilitar análise integrada de aspectos organizacionais e permitir o desenvolvimento de resposta concisa e clara.
Como nas demais questões, o desempenho dos estudantes foi fraco, evidenciando despreparo e falta de conhecimentos da Área.
Ficou clara a dificuldade de interpretação e de raciocínio lógico, além da frágil e imprecisa construção textual.
As respostas dadas pelos estudantes, na maioria, expressaram claramente uma assimilação deficiente, frágil e imprecisa, o que pode comprometer significativamente a prática desses futuros profissionais.
Ficou, também, evidente a dificuldade em analisar, interpretar o caso específico do contexto da questão e a associação deste à fundamentação teórica apropriada.
Mais uma vez, as respostas dadas denotam conteúdos memorizados e não apreendidos e, portanto, sem argumentação lógica própria. 3.3.3 Considerações Finais O conhecimento dos estudantes, de maneira geral, apresentou -se deficiente, frágil, impreciso ou incorreto.
Os estudantes demonstraram desempenho técnico e científico insuficiente e bastante deficiente ; baixa competência de estruturação lógica e analítica de pensamento, além de habilidade argumentativa praticamente inexistente.
Percebeu -se que o processo de ensino-aprendizagem parece privilegiar a memorização de conteúdos e não a construção propriamente dita de conhecimentos.

% matched lemmas: predominar
\end{textsample}
